\chapter{Extending the System}
\label{ch:extend}

The base-agent design makes some changes cheap and others expensive. This chapter walks through the common ones with working code.

\section{Adding a fourth specialist}

Say you want a Dinacharya agent that handles daily-routine and seasonal-regimen questions. Four edits.

\subsection{1. Tag the corpus}

Add a category in \code{AyurvedicTextProcessor.categorize\_content}. Position matters --- put it \emph{before} the \code{prakriti} test, or the dosha names in most routine passages will swallow it:

\begin{lstlisting}[style=py]
        # Daily and seasonal regimen - test before the dosha keywords
        if any(term in text_lower for term in
               ['dinacharya', 'ritucharya', 'daily routine', 'seasonal regimen',
                'sadvritta']):
            return 'regimen'

        # Prakriti-related content
        if any(term in text_lower for term in
               ['vata', 'pitta', 'kapha', 'constitution', 'prakriti']):
            return 'prakriti'
\end{lstlisting}

Re-run steps 1 and 2 after deleting \code{data/vectordb/}, since categories are baked into the index.

\subsection{2. Write the agent}

\begin{lstlisting}[style=py]
"""Regimen Agent - daily and seasonal routine"""
from .base_agent import BaseAgent
from typing import Optional

class RegimenAgent(BaseAgent):
    def __init__(self, rag_retriever, llm_client):
        super().__init__(name="Regimen Advisor", rag_retriever=rag_retriever,
                         llm_client=llm_client, temperature=0.3)

    def get_system_prompt(self) -> str:
        return """You are a senior Ayurvedic physician advising on Dinacharya
(daily regimen) and Ritucharya (seasonal regimen) for BAMS practitioners.

**CRITICAL INSTRUCTIONS**:
1. Ground every recommendation in the classical schedule, citing the Sutra
   Sthana chapter where the practice is described.
2. Give timings in both muhurta terms and clock time.
3. Use Sanskrit terms with brief English glosses.

**RESPONSE STRUCTURE**:
1. **Dinacharya Protocol**: ordered daily practices with timing
2. **Ritucharya Adjustment**: modifications for the current Ritu
3. **Contraindications**: practices to omit given the stated Vikriti
4. **Clinical Notes**: brief summary for a fellow practitioner"""

    def get_category_filter(self) -> Optional[str]:
        return "regimen"
\end{lstlisting}

\subsection{3. Teach the orchestrator about it}

Three touch points in \fpath{src/agents/orchestrator.py} --- the constructor, the router, and the cascade:

\begin{lstlisting}[style=py]
    def __init__(self, prakriti_agent, dosha_agent, treatment_agent,
                 regimen_agent, llm_client):
        ...
        self.regimen_agent = regimen_agent

    def analyze_query(self, query: str) -> Dict:
        ...
        needs_regimen = any(term in query_lower for term in
            ['routine', 'schedule', 'dinacharya', 'ritucharya', 'season',
             'wake', 'sleep time'])
        return {'prakriti': needs_prakriti, 'dosha': needs_dosha,
                'treatment': needs_treatment, 'regimen': needs_regimen}

    # in process_query, after the treatment block:
        if agent_activation['regimen']:
            regimen_info = base_additional_info.copy()
            if 'prakriti' in results:
                regimen_info['Prakriti'] = results['prakriti']
            if 'dosha' in results:
                regimen_info['Dosha Imbalance'] = results['dosha']
            regimen_result = self.regimen_agent.process(
                query, regimen_info, conversation_history)
            results['regimen'] = regimen_result['response']
            retrieved_sources.extend(regimen_result.get('sources', []))
\end{lstlisting}

Also add a branch in \code{synthesize\_response} so the new report reaches the final document:

\begin{lstlisting}[style=py]
        if 'regimen' in agent_results:
            synthesis_context += f"REGIMEN GUIDANCE:\n{agent_results['regimen']}\n\n"
\end{lstlisting}

\subsection{4. Wire it up}

In \fpath{src/ui/gradio\_app.py}:

\begin{lstlisting}[style=py]
        self.regimen_agent = RegimenAgent(self.retriever, self.llm_client)
        self.orchestrator  = OrchestratorAgent(
            self.prakriti_agent, self.dosha_agent, self.treatment_agent,
            self.regimen_agent, self.llm_client)
\end{lstlisting}

\begin{note}[title={The orchestrator is the friction point}]
Adding an agent costs one new file and four edits to one existing file. The activation dictionary, the cascade, and the synthesis context are all hand-written per agent. A registry --- a list of \code{(agent, trigger\_words, upstream\_keys)} tuples iterated in order --- would reduce this to one file, and is the refactor to reach for once a fifth agent appears.
\end{note}

\section{Adding a new LLM provider}

Implement two methods and the rest of the system does not care.

\begin{lstlisting}[style=py]
"""Anthropic client"""
import os
from typing import Optional, List, Dict
from anthropic import Anthropic

class AnthropicClient:
    def __init__(self, api_key: str = None, model: str = None):
        self.api_key = api_key or os.getenv("ANTHROPIC_API_KEY")
        if not self.api_key:
            raise ValueError("Anthropic API key not found. "
                             "Set ANTHROPIC_API_KEY in your .env file.")
        self.client     = Anthropic(api_key=self.api_key)
        self.model_name = model or os.getenv("ANTHROPIC_MODEL", "claude-sonnet-5")

    def generate(self, prompt: str, system_prompt: Optional[str] = None,
                 temperature: float = 0.5, max_tokens: int = 4000,
                 conversation_history: Optional[List[Dict]] = None, **kwargs) -> str:
        messages = []
        if conversation_history:
            # Drop the trailing turn: the UI already appended the current message
            for turn in conversation_history[:-1]:
                messages.append({"role": turn["role"], "content": turn["content"]})
        messages.append({"role": "user", "content": prompt})

        try:
            response = self.client.messages.create(
                model=self.model_name, system=system_prompt or "",
                messages=messages, temperature=temperature, max_tokens=max_tokens)
            return response.content[0].text
        except Exception as e:
            raise RuntimeError(f"Anthropic API Error: {e}")

    def generate_with_context(self, query: str, context: str, system_prompt: str,
                              temperature: float = 0.3, max_tokens: int = 4000,
                              conversation_history: Optional[List[Dict]] = None) -> str:
        prompt = f"""Context from Ayurvedic texts:

{context}

---

User Query: {query}

Based on the context provided above, please provide a response."""
        return self.generate(prompt=prompt, system_prompt=system_prompt,
                             temperature=temperature, max_tokens=max_tokens,
                             conversation_history=conversation_history)
\end{lstlisting}

Two details that matter. Expose \code{model\_name} --- \code{AyurMindApp} logs it on startup, and \code{OllamaClient} aliases \code{self.model\_name = self.model} for exactly this reason. And slice \code{conversation\_history[:-1]}, matching OpenRouter and Ollama rather than the Google and OpenAI clients (defect 3 in Chapter~\ref{ch:gaps}).

Then insert it into the fallback chain in \fpath{gradio\_app.py} at whatever priority you want.

\section{Adding a second source text}

The scraper is written against one MediaWiki instance, so a second text needs a second scraper class rather than a config entry. The processor and everything downstream are source-agnostic --- they only require the four metadata keys.

\begin{enumerate}[leftmargin=*]
  \item Write \code{SushrutaScraper} with the same \code{scrape\_all()} contract, writing into \code{data/raw/sushruta/} and producing a compatible \code{scraping\_summary.json}.
  \item Add a \code{text} key to chunk metadata (\code{"charaka"} / \code{"sushruta"}) so retrieval can be scoped per source, and extend \code{AyurvedicVectorStore.search} with a \code{text\_filter}.
  \item Re-run steps 1 and 2. The Chroma collection holds both corpora; the category filters keep working unchanged.
\end{enumerate}

The retrieval quality question is worth a thought before doing this. Chroma returns the nearest five vectors across the whole collection, so adding a second text does not give you five chunks from each --- it gives you five chunks total, and the two corpora compete. If per-source coverage matters, retrieve five per source with two filtered calls and concatenate.

\section{Making the source filter work}

The two-line fix for defect 2. In \code{AyurvedicTextProcessor.process\_all}, store the Sthana name alongside the display name:

\begin{lstlisting}[style=py]
                    chapter_metadata = {
                        'section':      section['section_name'],
                        'sthana':       section['section_key'].replace('_', ' '),
                        'section_code': section['section_code'],
                        'chapter':      chapter_info['title'],
                        ...
                    }
\end{lstlisting}

Then match against it in \code{AyurvedicVectorStore.search}:

\begin{lstlisting}[style=py]
        if source_filter:
            where_conditions.append({"$or": [{"sthana":  source_filter},
                                             {"section": source_filter},
                                             {"chapter": source_filter}]})
\end{lstlisting}

With \code{sthana} holding \code{"Siddhi Sthana"}, the phrase the orchestrator's regex extracts now matches exactly, and the treatment agent's per-Sthana prompt branches finally get the context they were written for.

\section{Things that are harder than they look}

\begin{table}[H]
\centering
\small
\begin{tabularx}{\textwidth}{@{}l X@{}}
\toprule
\textbf{Change} & \textbf{Why it is not small} \\
\midrule
Streaming responses & The cascade is inherently sequential and synthesis needs all specialist outputs before it starts. You can stream the final synthesis, but the first three-quarters of the wait produces nothing to show. Streaming per-agent progress into the chat is the realistic version. \\
\addlinespace
Hindi support & Needs a multilingual embedding model (so a rebuild of the whole index), translated system prompts, and a decision about whether retrieval happens in English with translated output or end-to-end in Hindi. The evaluation framework already defines the metric for it. \\
\addlinespace
Parallel agent execution & Dosha consumes Prakriti's output and Treatment consumes both, so only the Prakriti/Dosha pair is independent, and only when Treatment is inactive. The saving is one call on a subset of queries. \\
\addlinespace
Replacing keyword routing with a classifier & Straightforward to write --- one cheap model call returning a JSON activation dictionary --- but it adds a call to the fast path, which exists precisely to avoid extra calls. Worth measuring the misroute rate against the 40-case dataset's \code{expected\_agents} column first; that column exists for this purpose and nothing reads it yet. \\
\bottomrule
\end{tabularx}
\caption{Changes that touch the architecture rather than sitting on top of it.}
\end{table}
